\documentclass[11pt,a4paper]{article}
\usepackage[margin=22mm]{geometry}
\usepackage{amsmath,amssymb,booktabs,microtype}
\usepackage[hidelinks]{hyperref}
\newcommand{\dd}{\mathrm d}
\newcommand{\ii}{\mathrm i}
\newcommand{\tr}{\operatorname{tr}}
\title{Route G-R8: A Closed Guided Clock Protocol\\and Its Complete Effective Phase}
\author{Same spectral action, explicit controls, no fitted time law}
\date{29 September 2026}
\begin{document}
\maketitle
\begin{abstract}
Specify a symmetric guided interferometer, its full scalar potential,
preparation and analysis pulses. The ground propagation phase cancels
for this particular potential, whereas the internal-energy gravitational
phase survives. Both transport ramps are included. A finite-mass harmonic
diagnostic quantifies, rather than assumes away, the excited packet's
closure error. Pure/dephased controls and a partial-transpose witness
separate entanglement from visibility. The result is a closed conservative
effective protocol, not a demonstrated trapped-ion apparatus or an
all-orders relativistic/open-system prediction.
\end{abstract}

\section{Action and declared approximation}
Keep the G-R4--G-R7 scalar thermal gaps at common local temperature $T_0$:
\begin{equation}
 E_i=h\nu_i(T_0),\quad \nu_i(T_0)=\nu_i^0[1+b_i(T_0)],\quad
 H_C=\sum_{i=2,3}E_i|e_i\rangle\langle e_i| .
\end{equation}
Each run prepares $g$ and one selected $e_i$. These are not two independent
qubits in one ion. The ground mass $m$ absorbs any uniform ground-level
shift; no missing ground polarizability is assigned a numerical value.
Expand the same proper-time action in a weak uniform potential $gz$:
\begin{align}
 S_r&=\int\left[-mc^2+\tfrac12m\dot z^2-mgz-V_r
       +\ii\hbar\chi^\dagger\dot\chi
       -\chi^\dagger H_C\chi
          (1+gz/c^2-\dot z^2/(2c^2))\right]\dd t,\label{action}\\
 H_r&=mc^2+\frac{p^2}{2m}+mgz+V_r+
       H_C\left(1+\frac{gz}{c^2}-\frac{p^2}{2m^2c^2}\right).
       \label{ham}
\end{align}
The expansion keeps leading internal mass-energy coupling and Newtonian
external dynamics. It does not retain all post-Newtonian center-of-mass
terms. Metric, quantum mechanics and classical controls remain assumed.
Guided-clock formalism and compensated transport have established
precedents~\cite{roura,transport}; the following control card is our
explicit idealization, not an implementation imported from those papers.

There are two orthogonal guide modes $A,B$, initially/finally at the
same height. A unitary coupler mixes them into two output ports; two
orthogonal modes are never mapped into the same state. Other spatial
degrees of freedom are matched and factorized. Temperatures are locally
matched, not inferred from a global Tolman-equilibrium condition.
The time-dependent guide can supply work: a closed path does not mean
an energetically isolated atom. Controller backreaction and records are
excluded only by the stated classical-control approximation.

\newpage
\section{Full path and full potential, not just a force}
Let $t_r$ be one ramp duration, $t_h$ the hold and $t_f=2t_r+t_h$.
For $s(u)=10u^3-15u^4+6u^5$, define the arm separation
\begin{equation}
 q(t)=\begin{cases}
 d\,s(t/t_r),&0\leq t\leq t_r,\\
 d,&t_r<t<t_r+t_h,\\
 d[1-s((t-t_r-t_h)/t_r)],&t_r+t_h\leq t\leq t_f,
 \end{cases}\quad q_A=q/2,\quad q_B=-q/2 .\label{paths}
\end{equation}
Position, velocity and acceleration close. Specify the scalar potential
\begin{equation}
 V_r(z,t)=\frac{m\omega^2}{2}(z-q_r)^2-m(\ddot q_r+g)z+w_r(t),
 \qquad w_A=w_B=0 .\label{guide}
\end{equation}
Its uniform support field cancels the ground $mgz$ term.
At zeroth order in $H_C/(mc^2)$,
\begin{equation}
 \ddot z+\omega^2(z-q_r)=\ddot q_r .
\end{equation}
Thus the packet center follows $q_r$. The accelerating-frame wave
function is a rigid translated/boosted harmonic packet. For initial
oscillator eigenstate $n$, its phase is
\begin{equation}
 \psi_r(z,t)=
 e^{\ii[m\dot q_r(z-q_r)+S_r^c]/\hbar}
 e^{-\ii(\epsilon_n+mc^2)t/\hbar}\phi_n(z-q_r),\quad
 \dot S_r^c=\tfrac12m\dot q_r^2+m\ddot q_r q_r .\label{exactpacket}
\end{equation}
Direct substitution into the Schr\"odinger equation verifies this result.
Since $[q_r\dot q_r]_0^{t_f}=0$,
\begin{equation}
 S_r^c=-\frac m2\int_0^{t_f}\dot q_r^2\dd t,
 \qquad S_A^c-S_B^c=0 .\label{groundphase}
\end{equation}
The rest-energy and oscillator phases are common. At closure both
packets agree, up to the explicitly computed equal phases.

\paragraph{A physically important nonuniqueness.}
Replacing the last $z$ in Eq.~\eqref{guide} by $z-q_r$ changes no force
but adds $m(\ddot q_r+g)q_r$ to the scalar potential. Its arm difference
integrates to $mg\mathcal A$, where
\begin{equation}
 \mathcal A=\int_0^{t_f}q(t)\dd t=d(t_h+t_r),\qquad
 \int_0^{t_f}\dot q_r^2\dd t=\frac{5d^2}{7t_r}.\label{area}
\end{equation}
The resulting ground fringe changes by $-mg\mathcal A/\hbar$.
This is a different physical potential, not permission to discard a
branch phase. A compensating transformation of coupler phases would be
needed to interpret it as a mere basis convention. The specified
$w_A-w_B=0$ is a control requirement, not a consequence of force closure.

\newpage
\section{Pulses and the total observable phase}
At common height prepare the clock with $R_y(\pi/2)$, with its optical
phase defined by the local oscillator. Alternatively prepare its
energy-dephased mixture with the same populations. Apply the path pulse
\begin{equation}
 B(\alpha)=\exp[-\ii\pi(\cos\alpha\,X_P+\sin\alpha\,Y_P)/4]
           =\frac{I-\ii(\cos\alpha\,X_P+\sin\alpha\,Y_P)}{\sqrt2} .
\end{equation}
Use $B(\alpha_s)$ at $t=0$ and $B(\alpha_f)^\dagger$ after the full
return. The card sets $\alpha_s=0$ and scans $\alpha_f$. These are
impulse-limit operations, state-independent for the path coupler and
with negligible differential vertical clock recoil. Finite pulses and
real guide Stark shifts are not asserted to vanish in an apparatus.

Internal evolution follows directly from Eq.~\eqref{action}:
\begin{align}
 \theta_{ir}&=2\pi\nu_i(T_0)\int_0^{t_f}
      [1+gq_r/c^2-\dot q_r^2/(2c^2)]\dd t,\\
 \delta_i&=\theta_{iA}-\theta_{iB}
      =2\pi\nu_i(T_0)\frac{g\mathcal A}{c^2}.\label{delta}
\end{align}
A finite packet's common momentum variance also cancels from this
difference. The symmetric kinetic terms cancel; both ramps contribute
$dt_r$ to $\mathcal A$. During the stationary hold,
$\Delta\nu_i=\nu_i(T_0)gd/c^2$, consistent with the phase derivative.
For unequal baths one must instead integrate
$2\pi\nu_i^0[(1+b_{iA})\lambda_A-(1+b_{iB})\lambda_B]$ and account for
their forces and ground shifts. No such extra bath is silently introduced.

For $\rho_{gg}=\rho_{ee}=1/2$, direct multiplication of both path
pulses and the branch propagators gives
\begin{align}
 P(0,g)&=\tfrac14[1+\cos\Phi_g],&
 \Phi_g&=\alpha_f-\alpha_s,\\
 P(0,e_i)&=\tfrac14[1+\cos\Phi_{e_i}],&
 \Phi_{e_i}&=\alpha_f-\alpha_s-\delta_i .\label{port}
\end{align}
The other port uses the minus signs. Consequently
\begin{equation}
 P(0)=\tfrac12\{1+\cos(\delta_i/2)
             \cos[\alpha_f-\alpha_s-\delta_i/2]\},\quad
 \mathcal V=|\cos(\delta_i/2)|.\label{visibility}
\end{equation}
For this declared guide, the previously unresolved scalar phase is now
calculated, not simply omitted. Adding actual scalar offsets changes
both $\Phi_j$ by $-\int(w_A-w_B)\dd t/\hbar$. State-resolved phase
differences cancel a stable common offset; uncontrolled shot-to-shot
phase fluctuations can still destroy the signal.

All preparation/readout events occur at common height. Common clock
pulse and free-precession phases do not enter Eq.~\eqref{port}, but
must be calibrated for non-energy-basis joint measurements. There are
no unspecified laser kicks at the separated positions in this protocol.

\newpage
\section{Mass-dependent closure: a finite Gaussian diagnostic}
The same guide should not be declared exactly closing for every internal
mass. Test it with $m_i=m(1+\eta)$, $\eta=E_i/(mc^2)$ and
\begin{equation}
 H_{r,i}^{\rm diag}=m_i c^2+p^2/(2m_i)+m_i gz+V_r .\label{massdiag}
\end{equation}
This mass-only harmonic completion agrees with Eq.~\eqref{ham} to first
order in $\eta$, but does not supply omitted relativistic terms.
Write the actual center as $z_{A/B}=z_c\pm z_d/2$, with zero initial
position and velocity. Then $\Omega=\omega/\sqrt{1+\eta}$ and
\begin{align}
 \ddot z_c+\Omega^2z_c&=-\frac{\eta}{1+\eta}g,&
 z_c&=-\frac{\eta g}{\omega^2}(1-\cos\Omega t),\\
 z_d&=q+y,&
 \ddot y+\Omega^2y&=-\frac{\eta}{1+\eta}\ddot q .
\end{align}
Define $F(\Omega)=\int_0^{t_f}\ddot q(t)e^{\ii\Omega t}\dd t$.
The return acceleration is minus the outbound one, shifted by $t_r+t_h$:
\begin{equation}
 F(\Omega)=\frac d{t_r}
 \int_0^1s''(u)e^{\ii\Omega t_ru}\dd u\,
 [1-e^{\ii\Omega(t_r+t_h)}].\label{fourier}
\end{equation}
Let $I_c+\ii I_s=e^{\ii\Omega t_f}F(\Omega)^*$. At final closure of $q$,
\begin{equation}
 z_d(t_f)=-\frac{\eta I_s}{(1+\eta)\Omega},\quad
 \dot z_d(t_f)=-\frac{\eta I_c}{1+\eta},\quad
 \int_0^{t_f}z_d\dd t=\mathcal A+\frac{\eta I_c}{\omega^2}.
\end{equation}
Choose both $\omega t_r$ and $\omega t_h$ to be integer multiples of
$2\pi$. Then $F(\omega)=0$: first-order mass-dependent displacement
closes. Finite $\eta$ residuals are nonzero, not an exact all-mass theorem.

The exact phase of the two Gaussian packets' overlap includes the
endpoint term $-m_i\dot z_cz_d/\hbar$. The classical action difference is
$m_i z_c\dot z_d-\eta mg\int z_d\dd t$. Combining them yields
\begin{equation}
 \arg\langle\psi_{B,i}|\psi_{A,i}\rangle
 =-\delta_i-\frac{m\eta^2g}{\hbar\omega^2}\Re F(\Omega).
 \label{exactphase}
\end{equation}
Thus closing central positions alone would have missed a phase term.
For the initial ground Gaussian with covariance
$V_0=\operatorname{diag}(\hbar/(2m\omega),\hbar m\omega/2)$, set
\begin{equation}
 V=S_iV_0S_i^T,\quad
 S_i=\begin{pmatrix}\cos\Omega t_f&\sin\Omega t_f/(m_i\Omega)\\
 -m_i\Omega\sin\Omega t_f&\cos\Omega t_f\end{pmatrix}.
\end{equation}
Its overlap modulus is $r_i=e^{-D_i}$, where
\begin{equation}
 D_i=\frac{\Delta p^2V_{zz}+\Delta z^2V_{pp}
                  -2\Delta z\Delta pV_{zp}}{2\hbar^2},\quad
 \Delta z=z_d(t_f),\quad\Delta p=m_i\dot z_d(t_f).
\end{equation}
The energy-resolved excited fringe gets factor $r_i$, and the unconditioned
clock overlap becomes $(1+r_ie^{-\ii\delta_i^{\rm eff}})/2$.
This is a computed motional correction, not an invented environmental rate.

\newpage
\section{Pure/dephased comparison and a small-phase witness}
After known local phase changes, the ideal pure joint state is
\begin{equation}
 |\Psi_\delta\rangle=e^{\ii\delta Z_PZ_C/4}|++\rangle
 =\cos(\delta/4)|++\rangle+\ii\sin(\delta/4)|--\rangle .
\end{equation}
For example the diagonal internal evolution factors as
$e^{-\ii\bar\theta P_e}e^{-\ii\delta Z_P/4}
e^{\ii\delta Z_PZ_C/4}$, where $\bar\theta=(\theta_A+\theta_B)/2$.
The initial splitter's $-\ii$ phase is another local path rotation.
These phase changes require calibrated controls, not a claim that a
large optical phase is known from arbitrary decimal precision.

Replacing $|+\rangle\langle+|_C$ by $I_C/2$ leaves all energy-resolved
fringes and the reduced path state unchanged. The resulting joint state
is a convex sum of product states conditioned on $g,e$ and is separable.
For the pure state, $\mathcal N=|\sin(\delta/2)|/2$.
For the positive small phase in our card, use
\begin{equation}
 W=\tfrac14(I-X_PX_C+Z_PY_C+Y_PZ_C).\label{witness}
\end{equation}
Indeed $W=(|v_-\rangle\langle v_-|)^{T_C}$ with
$|v_-\rangle=(|+-\rangle-\ii|-+\rangle)/\sqrt2$ in the $X$ basis.
For any separable $\rho$, $\rho^{T_C}\geq0$, proving
$\tr(W\rho)\geq0$. Direct evaluation gives
\begin{equation}
 \langle W\rangle_{\rm pure}=-\tfrac12\sin(\delta/2),\qquad
 \langle W\rangle_{\rm dephased}=\tfrac14[1-\sin(\delta/2)] .
\end{equation}
For the opposite phase sign, reverse the signs of the two $Y$ terms.
This is an operational witness, not an inference from visibility.

\paragraph{Actual measurement settings.}
Measure $XX,ZY,YZ$ in separate runs. For either subsystem, analysis
$R_y(-\pi/2)$ followed by population ($Z$) detection measures $X$;
$R_x(\pi/2)$ followed by $Z$ measures $Y$; no pulse measures $Z$.
The three witness operators commute globally, despite requiring
incompatible local axes. To explicitly include noncommuting joint
controls, also measure $ZZ$:
\begin{equation}
 [Z_PY_C,Z_PZ_C]=2\ii I_PX_C\ne0 .
\end{equation}
Together with independently characterized coherent/dephased preparation,
these settings go beyond energy-basis populations. Finite-mass Gaussian
corrections, noise and pulse errors must be included in a measured state;
the ideal witness values are not observed violations.

\paragraph{Statistical cost should not be hidden.}
For the ideal pure state $\langle XX\rangle=1$ and
$\langle ZY\rangle=\langle YZ\rangle=-\sin(\delta/2)$.
With $N$ independent shots for each setting,
\begin{equation}
 \sigma_W^2=\frac{\cos^2(\delta/2)}{8N},\qquad
 N>\frac{9\cos^2(\delta/2)}{8\mathcal N^2}
\end{equation}
is the ideal three-standard-error separation criterion. This is not a
rigorous finite-sample confidence bound or an experimental sensitivity:
it excludes state-preparation/readout errors, drift, losses and count rate.

\newpage
\section{Numerical result, rejection criteria and next decision}
The single illustrative card fixes $d=1$ mm, $t_r=5$ ms, $t_h=10$ ms,
$\omega/(2\pi)=1$ kHz, $g=9.80665$ m/s$^2$, $T_0=300$ K and a rounded
$m=171\,u$. The latter is not a precision ionic-mass determination.
It gives $t_f=20$ ms, $\mathcal A=1.5\,10^{-5}$ m s and
$\Delta\tau=1.63670545\,10^{-21}$ s. Each arm reaches speed 0.1875 m/s
and acceleration 115.47 m/s$^2$; neither establishes hardware feasibility.
\begin{center}
\begin{tabular}{lrr}
\toprule
Quantity & E2 & E3\\
\midrule
Hold $\Delta\nu$ (Hz) & $7.51094\,10^{-5}$ & $7.00642\,10^{-5}$\\
Full $\delta$ (rad) & $7.07889\,10^{-6}$ & $6.60339\,10^{-6}$\\
Ideal $1-\mathcal V$ & $6.26384\,10^{-12}$ & $5.45060\,10^{-12}$\\
Ideal $\mathcal N=-\langle W\rangle$ & $1.76972\,10^{-6}$ & $1.65085\,10^{-6}$\\
Ideal shots per witness setting & $3.59205\,10^{11}$ & $4.12798\,10^{11}$\\
\bottomrule
\end{tabular}
\end{center}
The effective total phases are $\Phi_g=\alpha_f-\alpha_s$ and
$\Phi_e=\Phi_g-\delta$, not just a hold-only clock phase. Both pure
and dephased initial clocks have the same path visibility. The latter
has $\langle W\rangle\simeq0.25$, not a negative witness.

The harmonic diagnostic gives $|\Delta v|\simeq7.0\,10^{-23}$ m/s
(E2), $1-r\simeq5.2\,10^{-40}$ and a phase correction of order
$10^{-30}$ rad. These tiny numbers are cancellation checks in a perfect
model, not metrological claims. High precision evaluates the known
small-$\eta$ cancellation; it does not fit any parameter. Ordinary trap
errors will dominate. The modified centered potential of section2 would
instead add roughly $-3.96\,10^5$ rad of common ground phase while
leaving its forces unchanged.

\paragraph{Useful physical gates.}
An actual state-dependent guide energy $u_{i,r}(t)P_{e_i}$ adds
$\int(u_{i,A}-u_{i,B})\dd t/\hbar$ to $\delta_i$ and is degenerate
with the desired effect unless independently calibrated. For this card,
keeping that bias below the signal requires a time-averaged differential
frequency bias below roughly $5.6\,10^{-5}$ Hz (E2) or
$5.3\,10^{-5}$ Hz (E3), before demanding any precision margin.
Uncalibrated guide scalar offsets, controller which-path records, finite
pulse recoil, lifetime and radiation noise remain physical limitations.
The existing tens-of-ms E2 lifetime caution still applies to a 20 ms
sequence; no survival rate or complete dissipator is imported.

\textbf{Next decision:} choose and calibrate one real guide/coupler and
its state-resolved phase noise, or reject this implementation if the
requirements are not credible. Prioritize differential phase; pursue the
joint witness only with an adequate statistics/control budget. Do not
respond to poor feasibility by fitting a new time law, claiming exact
all-mass closure, or enlarging the derivation indefinitely.

\begin{thebibliography}{9}
\bibitem{roura} A. Roura, Phys. Rev. X10,021014(2020),
\href{https://doi.org/10.1103/PhysRevX.10.021014}{doi:10.1103/PhysRevX.10.021014}.
\bibitem{transport} E. Torrontegui et al., Phys. Rev. A83,013415(2011),
\href{https://doi.org/10.1103/PhysRevA.83.013415}{doi:10.1103/PhysRevA.83.013415}.
\bibitem{zych} M. Zych et al., Nat. Commun.2,505(2011),
\href{https://doi.org/10.1038/ncomms1498}{doi:10.1038/ncomms1498}.
\end{thebibliography}
\end{document}
